
on peut dans ce diagramme remplacer les \(Z^i\) par les \(A^i\), et on trouve un diagramme du type voulu, en termes des anneaux de Chow de \(X_t\) et \(X_s\).

\textbf{Remarque 7.13.13.} Soit \(T\) un schéma lisse et projectif sur un corps algébriquement clos. Il résulte trivialement du formalisme développé dans (SGA 5 IV) que si \(Z\) est un cycle de codimension \(i\) sur \(T\) dont l'image dans \(H^{2i}(T,\mathbb Q_\ell(i))\) est nulle, alors \(Z\) est numériquement équivalent à zéro. C'est une des conjectures standard de la théorie des cycles algébriques que la réciproque est également vraie. Il est immédiat d'ailleurs grâce à (VII 4.11) que l'on trouve une conjecture équivalente en y remplaçant le groupe des cycles algébriques par le groupe \(\Gr^\bullet(T)_{\mathbb Q}\), en utilisant la définition (7.13.8), ou encore par le groupe \(K^\bullet(T)\), en écrivant que l'équivalence numérique à zéro d'un élément de \(K^\bullet(T)\) est équivalente à la nullité de son augmentation et de ses classes de Chern dans les \(H^{2i}(T,\mathbb Q_\ell(i))\).

Ceci posé, il résulte immédiatement de la commutativité de (7.13.10) que si \(X_{\bar t}\) satisfait à la conjecture précédente, alors l'homomorphisme de spécialisation \(\Gr^\bullet(X_{\bar t})\longrightarrow \Gr^\bullet(X_{\bar s})\) (et a fortiori \(\Gr^\bullet(X_t)\longrightarrow \Gr^\bullet(X_s)\)) est compatible avec l'équivalence numérique, comme il a été conjecturé dans 7.9. On notera d'ailleurs que la conjecture à laquelle on a fait appel pour cela (équivalence numérique implique équivalence \(\mathbb Q_\ell\)-cohomologique) est plus faible que celle envisagée dans 7.12.5.

\subsection*{7.14. Homomorphisme de spécialisation pour les groupes de cycles et classes de cycles.}
Si \(T\) est un schéma noethérien, désignons par \(Z_i(T)\) (resp. \(Z^i(T)\)) le groupe des cycles algébriques sur \(T\) purement de dimension \(i\) (resp. purement de codimension \(i\)). Lorsque \(f:X\to S\) est de type fini, on définit alors
